\section{Related Literature}
\label{sec:literature}

Our paper connects to four strands of the literature on bid rigging and
procurement design.

\paragraph{Statistical tests for bid coordination.}
\citet{bajari2003deciding} develop two tests---exchangeability of bid
functions and conditional independence of bid residuals---to distinguish
competitive from collusive bidding. \citet{porter1993detection} and
\citet{porter1999ohio} detect bid rotation and market allocation in Ohio
school milk markets using structural estimation.
\citet{conley2016detecting} propose a method for detecting bidder groups
from co-bidding patterns without requiring bid values. Our FL
classification provides a natural partition of bidders into ``suspected
cover bidders'' (FL) and ``genuine competitors'' (non-FL) for applying
the Bajari--Ye framework, replacing the \emph{ad hoc} groupings
typically required.

\paragraph{Bid-rigging screens.}
\citet{imhof2019detecting} and \citet{imhof2018screening} develop
machine-learning screens using tender-level features (bid variance,
coefficient of variation, kurtosis) that achieve high accuracy against
known cartels. \citet{abrantesmetz2006a} propose variance screens;
\citet{hüschelrath2014cartel} review detection in procurement;
\citet{chassang2022robust} design robust screens for non-competitive
bidding. These approaches operate at the \emph{tender} level and require
bid values. Our FL screen complements them by operating at the
\emph{firm} level using only participation data, enabling a two-stage
detection workflow: FL screening first, then bid-level analysis of
flagged tenders.%
\footnote{Our comparison with Imhof-style screens % IJIO-R1: minor-5 — Imhof proxy caveat added
(Section~\ref{sec:robustness}) uses a CV median split as a coarser
approximation of the full \citet{imhof2019detecting} implementation,
which employs multiple bid-level features and a trained classifier. Our
horse-race therefore provides a conservative lower bound on the
incremental value of the FL screen relative to a fully implemented Imhof
classifier.}

\paragraph{Cover bidding.}
\citet{porter1999ohio} document complementary bidding where cartel
members submit deliberately high bids to create the appearance of
competition. \citet{pesendorfer2000study} analyzes school milk auctions
where incumbents coordinated bids. \citet{asker2010leniency} describes
the internal organization of a philatelic stamp cartel where members
rotated the role of cover bidder. \citet{marshall2012economics} provide a
comprehensive treatment of how cartels allocate cover bidding roles. Our
contribution is to identify cover bidders systematically from
participation data rather than \emph{ex post} from prosecuted cases.

\paragraph{Developing-country procurement.}
While the bid-rigging literature has focused predominantly on developed
economies, cartel enforcement is particularly challenging in developing
countries where competition authorities have limited resources and
procurement data systems are nascent
\citep{sanchezgraells2019screening}. Brazil's BEC platform provides a
unique laboratory: 4.5 million tender-items, 40 million bids, and 41,000
firms over 11 years, with both sealed-bid (convite) and electronic
auction (preg\~ao) modalities. The scale and institutional variation
enable analyses infeasible in smaller datasets.
